% =====================================================================
\chapter{Le Model Context Protocol dans ce projet}
\label{ch:mcp}

Ce chapitre répond point par point aux questions posées sur le MCP, avec le
code réel sous les yeux.

% ---------------------------------------------------------------------
\section{Qu'est-ce que le MCP Server dans \emph{ce} projet ?}

C'est un \textbf{processus Python séparé} qui expose les dix fonctions
d'exploitation Storage sous forme d'outils décrits, interrogeables par un modèle
de langage. Il n'est pas un serveur réseau : il ne écoute aucun port, il
dialogue par son entrée et sa sortie standard.

\begin{keybox}
\textbf{En une phrase.} Le MCP Server d'OASIS est le \emph{catalogue exécutable}
des capacités Storage : il dit au modèle quels outils existent, avec quels
paramètres, et il les exécute quand on le lui demande --- sans jamais décider
lui-même de ce qui est permis.
\end{keybox}

\noindent
La dernière partie de cette phrase est importante et le code l'assume
explicitement : \file{server.py} porte le commentaire « Applique par le persona
et la route Actions ; le serveur MCP reste neutre. » La politique de sécurité
vit \emph{au-dessus} du serveur, pas dedans.

% ---------------------------------------------------------------------
\section{Où se trouve son code ?}

\begin{longtable}{@{}p{4.6cm}p{2.0cm}p{7.8cm}@{}}
\toprule
\textbf{Fichier} & \textbf{Lignes} & \textbf{Rôle}\\
\midrule
\endhead
\file{app/mcp/storage/tool\_server.py} & 12 & Le serveur : instancie \cls{MCPServer} et enregistre les outils\\
\file{app/mcp/storage/server.py} & 60 & Le registre \fn{TOOL\_REGISTRY} et la classification du risque\\
\file{app/mcp/storage/client.py} & 118 & Le client \cls{StorageMCPClient} et le point d'audit\\
\file{app/mcp/storage/ssh.py} & 130 & L'accès SSH mutualisé par tous les outils\\
\file{app/mcp/storage/tools/*.py} & 7 fichiers & Les implémentations\\
\bottomrule
\end{longtable}

% ---------------------------------------------------------------------
\section{Comment est-il lancé ?}

Il n'est \textbf{jamais} lancé à la main. Le client le démarre comme
sous-processus, à la demande.

\begin{lstlisting}[language=Python,caption={app/mcp/storage/client.py --- paramètres de lancement}]
_BACKEND_DIR = Path(__file__).resolve().parents[3]      # -> backend/
_SERVER_PARAMS = StdioServerParameters(
    command=sys.executable,                              # le meme interpreteur Python
    args=["-m", "app.mcp.storage.tool_server"],          # module lance en -m
    cwd=str(_BACKEND_DIR),                               # pour que "app.*" soit importable
)
\end{lstlisting}

\begin{lstlisting}[language=Python,caption={app/mcp/storage/tool\_server.py --- le serveur complet}]
from mcp.server.mcpserver import MCPServer
from app.mcp.storage.server import TOOL_REGISTRY

mcp = MCPServer("oasis-storage-mcp")
for _tool in TOOL_REGISTRY.values():
    mcp.tool()(_tool)

if __name__ == "__main__":
    mcp.run(transport="stdio")
\end{lstlisting}

\subsection*{Cycle de vie côté client}

\begin{lstlisting}[language=Python,caption={Démarrage et arrêt}]
async def start(self) -> None:
    if self._session is not None:
        return
    self._stack = AsyncExitStack()
    read, write = await self._stack.enter_async_context(stdio_client(_SERVER_PARAMS))
    self._session = await self._stack.enter_async_context(ClientSession(read, write))
    await self._session.initialize()
\end{lstlisting}

\begin{keybox}
\textbf{Pourquoi aucune instance partagée au niveau du module.} Le commentaire
final de \file{client.py} l'explique : « une instance créée à l'import n'est
jamais démarrée, et un client non démarré échoue à chaque appel d'outil. Chaque
appelant ouvre le sien pour la durée de son travail. »

La route de confirmation ajoute l'argument économique : « Une confirmation est
un clic humain, donc rare : le coût du démarrage du sous-processus est
négligeable, et on évite le sous-processus partagé qui reste mort jusqu'au
prochain redémarrage du backend s'il tombe. »
\end{keybox}

% ---------------------------------------------------------------------
\section{Quels outils sont disponibles ?}

\begin{longtable}{@{}p{4.4cm}p{3.3cm}p{2.0cm}p{4.4cm}@{}}
\toprule
\textbf{Outil} & \textbf{Paramètres} & \textbf{Classe} & \textbf{Commande sur la VM}\\
\midrule
\endhead
\fn{get\_capacity} & \texttt{volume\_id?} & lecture & \texttt{df -T -B1 -P}\\
\fn{list\_backups} & \texttt{target?} & lecture & \texttt{systemctl list-timers} + \texttt{systemctl show}\\
\fn{get\_backup\_job\_status} & \texttt{target?} & lecture & alias de \fn{list\_backups}\\
\fn{get\_backup\_logs} & \texttt{job\_id}, \texttt{lines=20} & lecture & \texttt{journalctl -u ... -o short-iso}\\
\fn{list\_snapshots} & \texttt{volume\_id?} & lecture & \texttt{lvs -{}-noheadings ...}\\
\fn{get\_storage\_health} & --- & lecture & agrège capacité + sauvegardes\\
\fn{run\_backup} & \texttt{target} & \textbf{remédiation} & \texttt{systemctl start <service>}\\
\fn{create\_snapshot} & \texttt{volume\_id}, \texttt{volume\_group?}, \texttt{size\_gb=1.0} & \textbf{remédiation} & \texttt{sudo lvcreate -{}-snapshot}\\
\fn{repair\_backup\_repository} & \texttt{target}, \texttt{confirm=False} & \textbf{remédiation} & drop-in systemd + \texttt{daemon-reload}\\
\fn{restore\_from\_backup} & \texttt{target}, \texttt{backup\_id}, \texttt{confirm=False} & \textbf{destructif} & \texttt{sudo lvconvert -{}-merge}\\
\bottomrule
\end{longtable}

\begin{warnbox}
\textbf{\fn{sample\_disk\_performance} n'est pas un outil MCP.} Il vit dans
\file{tools/performance.py} mais n'est \emph{pas} dans le \fn{TOOL\_REGISTRY}.
Le modèle ne le voit donc jamais ; seul le collecteur de métriques l'appelle,
directement.
\end{warnbox}

% ---------------------------------------------------------------------
\section{Comment les outils sont-ils enregistrés ?}

En deux temps.

\subsection*{Temps 1 --- le registre}
\file{server.py} importe les fonctions et les place dans un dictionnaire :

\begin{lstlisting}[language=Python,caption={app/mcp/storage/server.py}]
TOOL_REGISTRY: dict[str, Callable[..., dict[str, Any]]] = {
    "get_capacity": get_capacity,
    "list_backups": list_backups,
    "run_backup": run_backup,
    "get_backup_job_status": get_backup_job_status,
    "get_backup_logs": get_backup_logs,
    "list_snapshots": list_snapshots,
    "create_snapshot": create_snapshot,
    "get_storage_health": get_storage_health,
    "repair_backup_repository": repair_backup_repository,
    "restore_from_backup": restore_from_backup,
}
\end{lstlisting}

\subsection*{Temps 2 --- le décorateur, appliqué en boucle}
\texttt{mcp.tool()(\_tool)} est l'application \emph{manuelle} du décorateur
\fn{@mcp.tool()}. Le SDK dérive alors, pour chaque fonction :
\begin{itemize}
  \item le \textbf{nom} de l'outil, depuis \texttt{\_\_name\_\_} ;
  \item la \textbf{description}, depuis le docstring ;
  \item le \textbf{schéma JSON des paramètres}, depuis les annotations de type
  et les valeurs par défaut.
\end{itemize}

\begin{keybox}
\textbf{Conséquence concrète et vérifiable.} Le docstring de
\fn{repair\_backup\_repository} est rédigé \emph{en anglais et à destination du
modèle} : il commence par « Repair a backup job whose restic repository cannot
be opened. Use this when \ldots{} the journal shows \texttt{unable to open
config file}. » Ce texte n'est pas de la documentation interne : c'est la
description que le LLM lit pour décider s'il doit appeler cet outil.
\end{keybox}

% ---------------------------------------------------------------------
\section{Comment le client communique-t-il avec le serveur ?}

\begin{figure}[h!]
\centering
\begin{tikzpicture}[node distance=8mm and 20mm]
\node[ambox,minimum width=34mm,minimum height=30mm] (cli) {\textbf{Processus backend}\\[2mm]\tiny uvicorn\\\tiny StoragePersona\\\tiny StorageMCPClient\\\tiny (asyncio.Lock)};
\node[greenbox,right=of cli,minimum width=34mm,minimum height=30mm] (srv) {\textbf{Sous-processus}\\[2mm]\tiny python -m app.mcp.storage.tool\_server\\\tiny MCPServer("oasis-storage-mcp")\\\tiny 10 outils enregistrés};
\node[redbox,right=of srv,minimum width=30mm,minimum height=22mm] (vm) {\textbf{VM Ubuntu}\\[2mm]\tiny sshd\\\tiny df / lvs / systemctl};

\draw[flowb] ([yshift=6mm]cli.east) -- node[above,lbl]{stdin : JSON-RPC} ([yshift=6mm]srv.west);
\draw[flowb] ([yshift=-6mm]srv.west) -- node[below,lbl]{stdout : résultat} ([yshift=-6mm]cli.east);
\draw[flowr] (srv) -- node[above,lbl]{paramiko / SSH 22} (vm);
\end{tikzpicture}
\caption{Le transport stdio : pas de socket, pas de port, un tube entre deux processus.}
\end{figure}

\subsection*{\fn{list\_tools\_for\_llm} --- traduction du catalogue}

\begin{lstlisting}[language=Python,caption={Du format MCP au format attendu par Groq}]
async def list_tools_for_llm(self) -> list[dict[str, Any]]:
    response = await self._require_session().list_tools()
    return [
        {"type": "function", "function": {
            "name": tool.name,
            "description": tool.description or tool.name,
            "parameters": getattr(tool, "inputSchema", None)
                          or getattr(tool, "input_schema", {}),
        }}
        for tool in response.tools
    ]
\end{lstlisting}

\noindent
Le double \fn{getattr} absorbe la différence de nommage entre versions du SDK
(\texttt{inputSchema} en camelCase, \texttt{input\_schema} en snake\_case).

\subsection*{\fn{call\_tool} --- l'appel et son audit}

\begin{enumerate}
  \item \textbf{Verrou} --- \texttt{async with self.\_lock} : une seule requête
  MCP à la fois sur une session partagée, le transport stdio n'étant pas
  multiplexé.
  \item \textbf{Appel} --- \fn{session.call\_tool(tool\_name, arguments)}.
  \item \textbf{Détection d'erreur} --- si \texttt{isError} (ou
  \texttt{is\_error}) est vrai, les blocs de texte sont concaténés et levés en
  \cls{RuntimeError}.
  \item \textbf{Décodage} --- le contenu textuel est passé à
  \texttt{json.loads} ; s'il est vide, on retombe sur
  \texttt{raw.model\_dump(mode="json")}.
  \item \textbf{Audit} --- si \texttt{db} et \texttt{user\_id} sont fournis,
  \fn{log\_mcp\_call} écrit dans \tbl{mcp\_calls}. Cet appel est lui-même
  entouré d'un \texttt{try/except} : « Un échec d'audit ne doit pas masquer le
  résultat de l'outil. »
  \item \textbf{Propagation} --- si le statut est \texttt{FAILED}, une
  \cls{RuntimeError} est levée \emph{après} l'écriture de l'audit.
\end{enumerate}

% ---------------------------------------------------------------------
\section{Comment une requête utilisateur arrive-t-elle jusqu'au MCP ?}

\begin{lstlisting}[caption={Trajet complet, fonction par fonction}]
1.  ChatPage.send()                       [frontend]
2.  api.post('/storage/chat', {...})      [frontend]
3.  chat()                                [api/v1/storage.py]
4.  storage_persona.handle_message()      [personas/storage/agent.py]
5.  async with StorageMCPClient() as mcp  [ouverture du sous-processus]
6.  mcp.list_tools_for_llm()              [catalogue -> format Groq]
7.  groq_client.chat_with_tools(...)      [le modele choisit un outil]
8.  mcp.call_tool(name, **arguments)      [client -> serveur]
9.  MCPServer dispatche vers la fonction  [sous-processus]
10. tools/capacity.get_capacity()         [implementation]
11. ssh.run("df -T -B1 -P")               [paramiko]
12. VM Ubuntu execute la commande
13. retour JSON -> mcp_calls -> messages  [audit + reinjection]
14. groq_client redige la reponse finale
15. ChatTurnOut -> React
\end{lstlisting}

% ---------------------------------------------------------------------
\section{Comment un outil est-il sélectionné ?}

\textbf{Par le modèle, et par lui seul}, à partir du catalogue reçu. Le code ne
contient aucun classifieur, aucune correspondance mot-clé \arrowto{} outil,
aucune heuristique de routage.

\begin{lstlisting}[language=Python,caption={La seule instruction de sélection}]
options.update({"tools": tools, "tool_choice": "auto"})
\end{lstlisting}

\noindent
Ce que le modèle voit pour décider : le prompt système, la date du serveur, le
contexte RAG, le bloc incident éventuel, l'historique, la question, et la
description de chaque outil. Ce que le code \emph{contraint} ensuite : ce qui
sera exécuté et ce qui sera seulement proposé.

% ---------------------------------------------------------------------
\section{Comment les paramètres sont-ils transmis ?}

\begin{enumerate}
  \item Le modèle renvoie \texttt{call.function.arguments}, une \textbf{chaîne
  JSON}.
  \item L'agent la décode dans un \texttt{try/except json.JSONDecodeError} et
  retombe sur \texttt{\{\}} en cas d'échec ; un résultat qui ne serait pas un
  dictionnaire est également ramené à \texttt{\{\}}.
  \item Les arguments sont passés en \texttt{**kwargs} à \fn{call\_tool}, puis
  transmis au serveur MCP, qui les applique à la fonction Python.
  \item Pour une action confirmée, les paramètres viennent de la colonne JSONB
  \texttt{action\_requests.parameters}, à laquelle
  \fn{confirm\_and\_execute\_action} ajoute \texttt{confirm=True}.
\end{enumerate}

\subsection*{Échappement avant la ligne de commande}
Tout argument qui finit dans une commande shell passe par \fn{ssh.quote}, un
alias de \texttt{shlex.quote}. Exemples réels :

\begin{lstlisting}[language=Python]
run(f"systemctl start {quote(service)}")
run(f"journalctl -u {quote(service)} -n {int(lines)} --no-pager -o short-iso")
run(f"sudo lvcreate --snapshot --size {float(size_gb)}G "
    f"--name {quote(snapshot_name)} {quote(source)}")
run(f"sudo lvconvert --merge {quote(backup_id)}")
\end{lstlisting}

\begin{keybox}
Les valeurs numériques sont converties (\texttt{int(lines)},
\texttt{float(size\_gb)}) plutôt que citées : une conversion ratée lève une
exception avant d'atteindre la VM. Les valeurs textuelles sont citées. Aucun
argument venu du modèle n'est concaténé brut dans une commande.
\end{keybox}

% ---------------------------------------------------------------------
\section{Comment le tool exécute-t-il l'action ?}

Prenons \fn{run\_backup} en entier :

\begin{lstlisting}[language=Python,caption={app/mcp/storage/tools/backup.py}]
def run_backup(target: str) -> dict:
    """
    Déclenche immédiatement un job de sauvegarde existant.
    Action non destructive (ajout de données), pas de confirmation requise.
    """
    service = _service_for(target) if target.endswith(".timer") else target
    try:
        run(f"systemctl start {quote(service)}")
    except StorageVMError as error:
        return {"error": str(error), "action": "backup_failed", "target": target}
    return {"action": "backup_started", "target": service, "status": "running"}
\end{lstlisting}

\noindent
Trois choses à remarquer :
\begin{itemize}
  \item la conversion \texttt{.timer} \arrowto{} \texttt{.service} est faite dans
  l'outil, le modèle n'a pas à la connaître ;
  \item l'erreur est \textbf{retournée}, pas levée --- c'est ce qui permet à
  \fn{confirm\_and\_execute\_action} de tester \texttt{result.get("error")} ;
  \item le docstring dit « pas de confirmation requise », mais
  \fn{run\_backup} appartient tout de même à \texttt{REMEDIATION\_TOOLS} : la
  politique de l'application est plus stricte que celle de l'outil.
\end{itemize}

\subsection*{La couche SSH}
\fn{ssh.run(command, allow\_failure=False)} : ouvre une connexion, exécute,
lit \texttt{stdout}, récupère le code de retour, lit \texttt{stderr}, ferme
dans un \texttt{finally}. Un code non nul lève \cls{StorageVMError}, sauf si
\texttt{allow\_failure=True} --- utile « aux commandes dont un code de retour non
nul est une réponse valide (\texttt{systemctl} sur une unit en échec, par
exemple) ».

\noindent
\fn{ssh.run\_many(commands)} exécute plusieurs commandes sur \emph{une seule}
connexion : « Ouvrir une session par commande coûte cher. »

% ---------------------------------------------------------------------
\section{Comment le résultat revient-il ?}

\begin{enumerate}
  \item L'outil renvoie un \texttt{dict} Python.
  \item Le SDK MCP le sérialise en JSON dans un bloc de contenu textuel.
  \item Le client concatène les blocs et exécute \texttt{json.loads}.
  \item Il l'écrit dans \tbl{mcp\_calls} (colonne \texttt{result}, JSONB).
  \item L'agent le passe à \fn{evaluate\_tool\_result} pour la détection
  d'anomalie.
  \item Il le réinjecte dans \texttt{messages} sous forme d'un message de rôle
  \texttt{tool}, avec le \texttt{tool\_call\_id} correspondant.
  \item Le modèle rédige la réponse au tour suivant.
\end{enumerate}

% ---------------------------------------------------------------------
\section{Comment les erreurs sont-elles gérées ?}

\begin{longtable}{@{}p{3.5cm}p{4.2cm}p{6.9cm}@{}}
\toprule
\textbf{Niveau} & \textbf{Erreur} & \textbf{Traitement}\\
\midrule
\endhead
SSH & Config incomplète & \cls{StorageVMError} : « SSH settings are incomplete in .env: \ldots »\\
SSH & VM injoignable & \cls{StorageVMError} : « Could not reach <host>: \ldots »\\
SSH & Commande en échec & \cls{StorageVMError} avec \texttt{stderr} ou le code de retour\\
Outil & \cls{StorageVMError} attrapée & \textbf{Retour} d'un dict \texttt{\{"error": ...\}}, jamais une exception\\
Client MCP & \texttt{isError} ou exception & \texttt{status=FAILED}, audit écrit, puis \cls{RuntimeError}\\
Agent & \cls{RuntimeError} & Convertie en \texttt{\{"status": "error", "message": ...\}} et réinjectée dans la conversation --- le modèle peut réagir\\
Route confirm & Serveur MCP non démarrable & HTTP 503, action laissée \texttt{PENDING} donc rejouable\\
Route confirm & Action introuvable ou déjà traitée & \cls{ValueError} \arrowto{} HTTP 404\\
Service & Résultat contenant \texttt{error} & Action marquée \texttt{FAILED}, pas \texttt{COMPLETED}\\
Audit & Échec d'écriture & \fn{logger.exception}, le résultat de l'outil est préservé\\
\bottomrule
\end{longtable}

\begin{keybox}
\textbf{La règle qui traverse toutes ces lignes.} Une erreur d'infrastructure ne
doit jamais devenir une donnée inventée, et ne doit jamais faire taire le
système. Elle remonte, nommée, jusqu'à l'écran.
\end{keybox}

% ---------------------------------------------------------------------
\section{Comment les permissions sont-elles gérées ?}

Quatre niveaux se superposent, et \textbf{aucun n'est dans le serveur MCP}.

\begin{longtable}{@{}p{2.2cm}p{4.0cm}p{8.4cm}@{}}
\toprule
\textbf{Niveau} & \textbf{Où} & \textbf{Ce qui est contrôlé}\\
\midrule
\endhead
1 & \fn{get\_current\_user} & Jeton valide, compte actif\\
2 & \fn{require\_role} & Rôle autorisé sur la route\\
3 & \texttt{REMEDIATION\_TOOLS} + \fn{\_propose\_action} & Le chat ne peut pas exécuter un outil qui modifie la VM\\
4 & \texttt{confirm=True} dans les outils & \fn{restore\_from\_backup} et \fn{repair\_backup\_repository} refusent d'agir sans ce drapeau, seul \route{/confirm} l'accorde\\
\bottomrule
\end{longtable}

\begin{warnbox}
\textbf{Deux constats de lecture, à énoncer tels quels.}

\emph{(a)} \texttt{DESTRUCTIVE\_TOOLS} et la fonction \fn{risk\_for} sont
définis dans \file{server.py} mais \textbf{ne sont appelés nulle part} dans le
dépôt. La qualification « HIGH / LOW » n'est donc pas écrite en base
aujourd'hui. Le frontend possède sa \emph{propre} liste
(\texttt{DESTRUCTIVE\_ACTIONS = new Set(['restore\_from\_backup'])}) pour
afficher l'avertissement rouge --- une duplication à signaler.

\emph{(b)} La fonction \fn{invoke} de \file{server.py} n'est pas utilisée non
plus : le serveur MCP enregistre les fonctions directement, il ne passe pas par
elle.
\end{warnbox}

% ---------------------------------------------------------------------
\section{L'architecture LLM \arrowto{} Agent \arrowto{} MCP Client \arrowto{} MCP Server \arrowto{} Tool \arrowto{} Infrastructure correspond-elle au code ?}

\textbf{Oui, avec deux précisions.}

\subsection*{Correspondance exacte}

\begin{longtable}{@{}p{3.4cm}p{4.6cm}p{6.6cm}@{}}
\toprule
\textbf{Couche théorique} & \textbf{Réalisation dans le code} & \textbf{Fichier}\\
\midrule
\endhead
LLM & Groq, \texttt{openai/gpt-oss-120b} & \file{llm/groq\_client.py}\\
Agent & \cls{StoragePersona} & \file{personas/storage/agent.py}\\
MCP Client & \cls{StorageMCPClient} & \file{mcp/storage/client.py}\\
MCP Server & \cls{MCPServer("oasis-storage-mcp")} & \file{mcp/storage/tool\_server.py}\\
Tool & 10 fonctions Python & \file{mcp/storage/tools/}\\
Infrastructure & VM Ubuntu par SSH & \file{mcp/storage/ssh.py}\\
\bottomrule
\end{longtable}

\subsection*{Précision 1 --- l'agent n'est pas un simple relais}
Entre l'Agent et le MCP Client s'intercale un \textbf{garde-fou} : quatre des
dix outils ne traversent jamais cette frontière depuis le chat. Le schéma
générique laisserait croire que tout appel du modèle atteint l'infrastructure.
Ici, ce n'est vrai que pour les six outils de lecture.

\subsection*{Précision 2 --- une voie parallèle contourne MCP}
Le dashboard, la surveillance et le collecteur appellent les mêmes fonctions
\emph{sans} passer par le client ni le serveur MCP. La chaîne complète n'est
donc pas la seule voie d'accès à l'infrastructure --- c'est la voie de
\emph{l'agent}.

\begin{figure}[h!]
\centering
\begin{tikzpicture}[node distance=6mm and 14mm]
\node[vibox,minimum width=26mm] (llm) {LLM (Groq)};
\node[vibox,right=of llm,minimum width=26mm] (ag) {Agent\\\tiny StoragePersona};
\node[ambox,right=of ag,minimum width=26mm] (cl) {MCP Client};
\node[ambox,right=of cl,minimum width=26mm] (sv) {MCP Server};
\node[greenbox,below=14mm of sv,minimum width=26mm] (to) {Tool Python};
\node[redbox,below=of to,minimum width=26mm] (in) {Infrastructure\\\tiny VM Ubuntu};
\node[bluebox,left=of to,minimum width=30mm] (direct) {Dashboard · Watcher\\Collector\\\tiny appel direct};

\draw[flow,<->] (llm)--(ag);
\draw[flow] (ag)--(cl); \draw[flow] (cl)--(sv); \draw[flow] (sv)--(to);
\draw[flowr] (to)--(in);
\draw[flowb] (direct)--(to);
\node[lbl,below=1mm of ag] {garde-fou :\\4 outils arrêtés ici};
\end{tikzpicture}
\caption{Les deux voies d'accès aux outils.}
\end{figure}
